\chapter*{Conclusion}
This paper proposes a Quantum Retrieval-Augmented Generation (QRAG) framework, which combines quantum computing and classical RAG models to improve syntactic ambiguity resolution by RAG systems. The paper shows that classical RAG systems function well in processing linear and keyword queries; however, their performance is negatively impacted by syntactic ambiguity because conventional parsing methods cannot handle it adequately.

The proposed QRAG system employs Quantum Natural Language Processing (QNLP) to solve this problem. An Ambiguity Controller component determines whether quantum resources should be used based on the structural complexity of the input query. Thus, the hybrid approach helps to maximize efficiency and improve the quality of syntactic ambiguity resolution.

The experimental results suggest that quantum processing improves syntactic ambiguity resolution significantly compared to a pure classical approach. The paper proves the concept of the Utility Threshold, meaning that quantum computation becomes more beneficial than classical processing only after a certain level of syntactic complexity. Therefore, using hybrid processing allows us to combine the advantages of both approaches.

Importantly, this research reveals valuable information regarding the limits of classical RAG systems' effectiveness. For example, a linearity trap occurs when there is a high lexical similarity between queries but no semantic similarity. This flaw makes classical systems inefficient because they are unable to adjust the context. Therefore, contextual realignment is essential for producing accurate and relevant responses.

To conclude, quantum computation can substantially augment the performance of RAG systems. However, its application should be selective since classical RAG systems do not need quantum resources in processing simple queries.